\chapter[Data Representation and Memory Model]{Data Representation\\and Memory Model}
\label{ch:data-representation}

\section{Overview}

This chapter defines the representation rules used by registers, memory, program images, exception vectors, and
multi-word values. These rules are architectural: assemblers, loaders, emulators, hardware implementations, and
debuggers must agree on them.

\section{Word and Byte Ordering}

The SIRC-1 is a word-addressed processor. Each architectural memory address names one 16-bit word. Byte addressing is
not provided by the CPU instruction set.

When a word is represented outside the CPU as bytes, such as in a ROM image or debugger dump, the word is stored in
big-endian byte order: the most significant byte appears first.

\begin{table}[H]
	\centering
	\begin{tabular}{lll}
		\toprule
		\textbf{Value}  & \textbf{First byte} & \textbf{Second byte} \\
		\midrule
		\texttt{0x1234} & \texttt{0x12}       & \texttt{0x34}        \\
		\texttt{0xFF80} & \texttt{0xFF}       & \texttt{0x80}        \\
		\bottomrule
	\end{tabular}
	\caption{External Byte Order for 16-bit Words}
	\label{tab:external-byte-order}
\end{table}

Multi-word quantities are stored with the most significant word at the lowest word address. This applies to
instructions, exception vectors, saved addresses, and any software-defined multi-word integer convention that chooses
to use the architectural order.

\section{Address Values}

Architectural addresses are 24-bit word addresses. The four address-register pairs (\reg{l}, \reg{a}, \reg{s}, and
\reg{p}) represent an address by combining the low byte of the high register with the full low register:

\begin{equation}
	\text{Address} = ((\text{HighRegister} \mathbin{\&} \texttt{0x00FF}) \ll 16) \mid \text{LowRegister}
\end{equation}

Bits 15--8 of the high register are not driven on the address bus and do not participate in address comparisons,
effective-address calculation, instruction fetch, or exception vector dispatch. They are preserved as ordinary register
storage unless an instruction explicitly writes the high register.

\begin{table}[H]
	\centering
	\begin{tabular}{llll}
		\toprule
		\textbf{Pair value} & \textbf{High register} & \textbf{Low register} & \textbf{Address bus value} \\
		\midrule
		\reg{ah:al}         & \texttt{0x0012}        & \texttt{0x3456}       & \texttt{0x123456}          \\
		\reg{ah:al}         & \texttt{0xAB12}        & \texttt{0x3456}       & \texttt{0x123456}          \\
		\reg{ph:pl}         & \texttt{0x0000}        & \texttt{0x0100}       & \texttt{0x000100}          \\
		\bottomrule
	\end{tabular}
	\caption{Address Register Pair Interpretation}
	\label{tab:address-pair-interpretation}
\end{table}

When a 24-bit address is stored in memory as a 32-bit quantity, the high word is stored first and the low word second.
Only bits 7--0 of the high word are significant when the value is loaded as an address.

\begin{table}[H]
	\centering
	\begin{tabular}{lll}
		\toprule
		\textbf{Stored address} & \textbf{Word at address N} & \textbf{Word at address N+1} \\
		\midrule
		\texttt{0x123456}       & \texttt{0x0012}            & \texttt{0x3456}              \\
		\texttt{0xFECAFE}       & \texttt{0x00FE}            & \texttt{0xCAFE}              \\
		\bottomrule
	\end{tabular}
	\caption{Stored 24-bit Address Format}
	\label{tab:stored-address-format}
\end{table}

\section{Instruction Storage}

Every instruction is 32 bits, stored as two consecutive 16-bit words:

\begin{itemize}
	\item Word at \reg{p}: instruction bits 31--16
	\item Word at \reg{p}+1: instruction bits 15--0
\end{itemize}

Instruction fetch reads the high word first, increments \reg{p} by one word, reads the low word, then increments
\reg{p} by one word again. Instruction addresses must be even word addresses. Fetching from an odd word address raises
an alignment fault.

\section{Exception Vector Storage}
\label{sec:exception-vector-storage}

Each exception vector table entry is a stored 24-bit address in the format shown in
Table~\ref{tab:stored-address-format}. Vector entry \texttt{V} starts at word address \texttt{V * 2}:

\begin{itemize}
	\item Word at \texttt{V * 2}: target address high word
	\item Word at \texttt{V * 2 + 1}: target address low word
\end{itemize}

For example, vector \texttt{0x06} occupies word addresses \texttt{0x00000C} and \texttt{0x00000D}. If it should jump to
\texttt{0x123456}, those two words contain \texttt{0x0012} and \texttt{0x3456}.

\section{Immediate Values}

Immediate fields are encoded as raw two's-complement bit patterns. The consuming instruction defines whether the field
is interpreted as signed, unsigned, or bitwise data.

\begin{table}[H]
	\centering
	\begin{tabularx}{\textwidth}{llX}
		\toprule
		\textbf{Use}                   & \textbf{Width} & \textbf{Interpretation}                                              \\
		\midrule
		ALU arithmetic immediate       & 16 bits        & Two's-complement operand; result is truncated to 16 bits             \\
		ALU arithmetic short immediate & 8 bits         & Zero-extended operand in the range \texttt{0x00}--\texttt{0xFF}      \\
		ALU logical immediates         & 16 or 8 bits   & Bit mask; no sign interpretation is applied by the logical operation \\
		Memory displacements           & 16 bits        & Two's-complement offset added to the low address word                \\
		Coprocessor operands           & 16 bits        & Coprocessor-defined command or parameter bits                        \\
		\bottomrule
	\end{tabularx}
	\caption{Immediate Field Interpretation}
	\label{tab:immediate-interpretation}
\end{table}

Short immediates occupy an 8-bit field and are zero-extended before use by ALU arithmetic. They are therefore suitable
for compact non-negative constants and bit masks, not for encoding negative arithmetic operands. Assemblers must reject
values that cannot be represented in the selected immediate field width unless they deliberately offer a documented
truncation mode.

Short immediates are not used for memory, effective-address, branch, jump, or subroutine-call displacements. Those
forms use 16-bit immediate displacements, register displacements, or assembler-resolved full-width branch offsets.

\section{Arithmetic and Address Wraparound}

General-purpose register arithmetic is performed modulo $2^{16}$. Arithmetic condition flags describe the truncated
16-bit result: C records unsigned carry or borrow, and V records signed overflow.

Address calculation for memory, effective-address, branch, jump, and auto-update operations adds offsets to the low
word of the selected address register pair. The high address-register word is not changed by ordinary displacement,
pre-decrement, or post-increment calculation. If the low word overflows or underflows and
\statusbit{TrapOnAddressOverflow} is set, the instruction raises a segment-overflow fault. If the bit is clear, the low
word wraps modulo $2^{16}$.

The program counter follows the same low-word wrap rule during normal instruction fetch. Since instructions are two
words long and must begin at even word addresses, a fetch that would begin from an odd \reg{pl} value raises an
alignment fault.

\section{Memory Ordering}

The SIRC-1 has no architectural cache, store buffer, memory reordering, or speculative memory access. Bus reads and
writes issued by the CPU are externally visible in program order. A bus device may delay completion by withholding bus
acknowledgement, but this stalls the current execution phase rather than allowing later bus accesses to pass it.
